\documentclass[aspectratio=169,12pt]{beamer}
\usepackage{amsmath}
\usepackage{amssymb}
\usepackage{array}
\usepackage[most]{tcolorbox}
\usepackage{graphicx}
\usepackage[sfdefault,lf]{FiraSans}
\renewcommand{\familydefault}{\sfdefault}
\setlength{\parindent}{0pt}
\everymath{\displaystyle}

\setbeamertemplate{navigation symbols}{}
\setbeamertemplate{footline}{}
\setbeamertemplate{headline}{}
\setbeamertemplate{blocks}[rounded][shadow=false]

\definecolor{mmpageWhite}{HTML}{FCFBF7}
\definecolor{mmtanSoft}{HTML}{F7F5EF}
\definecolor{mmgreenDeep}{HTML}{244B2C}
\definecolor{mmgreenLeaf}{HTML}{5D8A56}
\definecolor{mmgreenSoft}{HTML}{E4EFD9}
\definecolor{mmtext}{HTML}{163221}
\definecolor{mmwarn}{HTML}{8F2F36}
\definecolor{mmwarnSoft}{HTML}{F8E8EA}

\setbeamercolor{background canvas}{bg=mmpageWhite}
\setbeamercolor{normal text}{fg=mmtext,bg=mmpageWhite}
\setbeamercolor{block title}{fg=mmgreenDeep,bg=mmgreenSoft}
\setbeamercolor{block body}{fg=mmtext,bg=mmtanSoft}
\setbeamercolor{block title example}{fg=mmgreenDeep,bg=mmgreenSoft}
\setbeamercolor{block body example}{fg=mmtext,bg=white}
\setbeamercolor{block title alerted}{fg=mmwarn,bg=mmwarnSoft}
\setbeamercolor{block body alerted}{fg=mmtext,bg=white}
\setbeamercolor{itemize item}{fg=mmgreenDeep}
\setbeamercolor{itemize subitem}{fg=mmgreenDeep}

\newtcolorbox{MMCard}[2][]{enhanced,breakable=false,colback=#2,colframe=black,boxrule=1.5pt,arc=8pt,left=5pt,right=5pt,top=4pt,bottom=4pt,before skip=0pt,after skip=0pt,#1}
\newcommand{\KoruIcon}{\raisebox{-0.2em}{\includegraphics[height=1.05em]{../../../manamaths/SITE/header-logo.png}}}
\newcommand{\NotesTitle}[1]{{\colorbox{mmtanSoft}{\parbox{0.975\linewidth}{\vspace{0.14em}\hspace{0.25em}{\LARGE\bfseries\textcolor{mmgreenDeep}{#1}\hfill\KoruIcon}\vspace{0.14em}}}\par\vspace{0.18em}{\color{mmgreenLeaf}\rule{\linewidth}{2.0pt}}\par\vspace{0.12em}}}
\newcommand{\SectionTitle}[1]{{\large\bfseries\textcolor{mmgreenDeep}{#1}}\par\vspace{0.04em}}
\newcommand{\GreenDivider}{\par\vspace{0.01em}{\color{mmgreenLeaf}\rule{\linewidth}{2.4pt}}\par\vspace{0.03em}}
\newcommand{\KeyIdea}[1]{\begin{MMCard}[title={\textbf{Key idea}},colbacktitle=mmgreenSoft,coltitle=mmgreenDeep]{mmtanSoft}\small #1\end{MMCard}}
\newcommand{\WorkedExample}[2]{\begin{MMCard}[title={\textbf{#1}},colbacktitle=mmgreenSoft,coltitle=mmgreenDeep]{white}\small #2\end{MMCard}}
\newcommand{\CommonMistake}[1]{\begin{MMCard}[title={\textbf{Common mistake}},colbacktitle=mmwarnSoft,coltitle=mmwarn]{white}\small #1\end{MMCard}}
\newcommand{\TryThese}[1]{\begin{MMCard}[title={\textbf{Try these}},colbacktitle=mmgreenSoft,coltitle=mmgreenDeep]{mmtanSoft}\small #1\end{MMCard}}
\newcommand{\StepLine}[2]{\textbf{#1.} #2\par\vspace{0.03em}}

\usepackage{tikz}

\setbeamertemplate{background}{
 \begin{tikzpicture}[remember picture,overlay]
 \draw[
 black,
 line width=1.5pt,
 rounded corners=10pt
 ]
 ([xshift=0.45cm,yshift=-0.45cm]current page.north west)
 rectangle
 ([xshift=-0.45cm,yshift=0.45cm]current page.south east);
 \end{tikzpicture}
}

\begin{document}

\begin{frame}[t,shrink=8]
\NotesTitle{Notes \& Steps}
\footnotesize
\begin{columns}[T,onlytextwidth]
\begin{column}{0.48\textwidth}
\KeyIdea{Scientific notation writes very large or very small numbers compactly as \(a \times 10^{n}\), where \(1 \le a < 10\) and \(n\) is an integer.}
\SectionTitle{Steps — large numbers}
\StepLine{1}{Move the decimal point left until only one digit remains before it.}
\StepLine{2}{Count how many places you moved — that is the exponent \(n\).}
\StepLine{3}{Write the result as \(a \times 10^{n}\).}
\SectionTitle{Steps — small numbers}
\StepLine{1}{Move the decimal point right until only one non-zero digit is before it.}
\StepLine{2}{Count the places moved — that is \(-n\).}
\StepLine{3}{Write as \(a \times 10^{-n}\).}
\end{column}
\begin{column}{0.48\textwidth}
\SectionTitle{Examples}
\begin{itemize}
  \setlength{\itemsep}{0.12em}
  \item \(2000 = 2 \times 10^{3}\) (move 3 left)
  \item \(0.00037 = 3.7 \times 10^{-4}\) (move 4 right)
  \item \(5 \times 10^{3} = 5000\) (move 3 right)
  \item \(7.5 \times 10^{-2} = 0.075\) (move 2 left)
\end{itemize}
\end{column}
\end{columns}
\GreenDivider
\vfill
\begin{columns}[b,onlytextwidth]
\begin{column}{0.48\textwidth}
\CommonMistake{Writing \(30 \times 10^{2}\) instead of \(3 \times 10^{3}\). The coefficient must be between 1 and 10. \(30\) is too large. Adjust: \(30 \times 10^{2} = 3 \times 10^{3}\).}
\end{column}
\begin{column}{0.48\textwidth}
\TryThese{\begin{enumerate}
  \item Write \(2000\) in scientific notation.
  \item Write \(0.00037\) in scientific notation.
  \item Express \(5 \times 10^{3}\) as an ordinary number.
\end{enumerate}}
\end{column}
\end{columns}
\end{frame}

\begin{frame}[t,shrink=12]
\NotesTitle{Notes \& Steps}
\begin{columns}[T,onlytextwidth]
\begin{column}{0.48\textwidth}
\WorkedExample{Example 1: large number}{Write \(9,500,000\) in scientific notation.\[9\,500\,000\]Move the decimal point 6 places left: \(9.5\). Answer: \(9.5 \times 10^{6}\).}
\end{column}
\begin{column}{0.48\textwidth}
\WorkedExample{Example 2: small number}{Write \(0.00012\) in scientific notation.\[0.00012\]Move the decimal point 4 places right: \(1.2\). Answer: \(1.2 \times 10^{-4}\).}
\end{column}
\end{columns}
\GreenDivider
\TryThese{\begin{enumerate}
  \item Convert \(3.84 \times 10^{5}\) to a normal number.
  \item Write \(0.02\) in scientific notation.
  \item Write \(7.5 \times 10^{-2}\) as a decimal.
\end{enumerate}}
\end{frame}

\begin{frame}[t,shrink=12]
\NotesTitle{Notes \& Steps}
\begin{columns}[T,onlytextwidth]
\begin{column}{0.48\textwidth}
\WorkedExample{Example 3: multiplying in sci. not.}{Calculate \((3 \times 10^{4}) \times (2 \times 10^{3})\).\[3 \times 2 = 6\]\[10^{4} \times 10^{3} = 10^{7}\]Answer: \(6 \times 10^{7}\).}
\end{column}
\begin{column}{0.48\textwidth}
\WorkedExample{Example 4: dividing in sci. not.}{Calculate \((6 \times 10^{5}) \div (2 \times 10^{2})\).\[6 \div 2 = 3\]\[10^{5} \div 10^{2} = 10^{3}\]Answer: \(3 \times 10^{3}\).}
\end{column}
\end{columns}
\GreenDivider
\CommonMistake{Adding exponents instead of subtracting during division. Remember: divide coefficients, subtract exponents.}
\end{frame}

\end{document}
